% PAGES: 159-160
% Theorem 5.6's proof opens here (page 160 / folio 144) and is left OPEN at
% the end of this file: the last line is equation (5.22), immediately after
% which the source continues "where Q is an orthogonal matrix; cf. Theorem
% 5.4." at the top of page 161 -- that is the first line of sec05_c1_4.tex.
% Do not print a second "\begin{proof}" there.

\begin{lemma}
Let $M$ be an $m \times n$ matrix.

(i) The matrix $M\tp M$ is an $n \times n$ symmetric positive semidefinite matrix.

(ii) The matrix $M\tp M$ is symmetric positive definite if and only if the columns of
$M$ are linearly independent.\footnote{A necessary (but not sufficient) condition for
the $m \times n$ matrix $M$ to have linearly independent columns is that it cannot
have more columns than rows, i.e., $n \leq m$.}
\end{lemma}

\begin{proof}
Note that the matrix $M\tp M$ is symmetric, since
\[
(M\tp M)\tp \;=\; M\tp(M\tp)\tp \;=\; M\tp M.
\]

(i) Let $x = (x_i)_{i=1:n}$ be a column vector in $\R^n$. Then,
\begin{equation}
x\tp M\tp Mx \;=\; (Mx)\tp Mx \;=\; (Mx,Mx) \;=\; \|Mx\|^2.
\end{equation}
Since $\|Mx\|^2 \geq 0$, it follows from (5.20) that
\begin{equation}
x\tp M\tp Mx \;\geq\; 0, \quad \forall\, x \in \R^n,
\end{equation}
and therefore that the matrix $M\tp M$ is symmetric positive semidefinite; cf.
Definition 5.6.

(ii) Let $M = \mathrm{col}\,(m_k)_{k=1:n}$. From (5.20), it follows that
\[
x\tp M\tp Mx = 0 \iff \|Mx\|^2 = 0 \iff Mx = 0 \iff \sum_{k=1}^{n} x_km_k = 0,
\]
where the matrix-vector multiplication formula (1.7) was used for the last
equivalence.

Thus, there exists a vector $x \neq 0$ such that $x\tp M\tp Mx = 0$ if and only if the
columns of the matrix $M$ are not linearly independent. Since $x\tp M\tp Mx \geq 0$
for all $x \in \R^n$, see (5.21), we conclude from Definition 5.6 that the matrix
$M\tp M$ is symmetric positive definite if and only if the columns of $M$ are linearly
independent.
\end{proof}

% "of the from" is the source's own wording (an apparent typo for "of the
% form") -- reproduced as printed, per the never-correct-the-author rule.
Note that symmetric positive semidefinite matrices of the from $M\tp M$ appear often
in practical applications such as finding covariance matrices from time series data,
solving least square problems, and doing linear regression; see Chapter 7 and
Chapter 8.

\par\medskip\noindent\textit{Example:}\ Show that the $N \times N$ matrix
\[
B_N \;=\; \begin{pmatrix}
2 & -1 & \dots & 0\\
-1 & \ddots & \ddots & \vdots\\
\vdots & \ddots & \ddots & -1\\
0 & \dots & -1 & 2
\end{pmatrix}
\]
is symmetric positive definite.\footnote{Another proof of this fact can be given using
the fact that a matrix is symmetric positive definite if and only if all its
eigenvalues are positive; see Theorem 5.6 and the example thereafter.}

\par\medskip\noindent\textit{Solution:}\ Recall from (10.20) that
\[
x\tp B_Nx \;=\; \sum_{1 \leq j,k \leq N} B_N(j,k)x_jx_k,
\]
and therefore
\begin{align*}
x\tp B_Nx &\;=\; \sum_{i=1}^{N} 2x_i^2 \;-\; \sum_{i=2}^{N} x_{i-1}x_i \;-\;
\sum_{i=1}^{N-1} x_ix_{i+1}\\
&\;=\; 2\sum_{i=1}^{N} x_i^2 \;-\; 2\sum_{i=1}^{N-1} x_ix_{i+1}\\
&\;=\; x_1^2 \;+\; \sum_{i=1}^{N-1} x_{i+1}^2 \;+\; \sum_{i=1}^{N-1} x_i^2 \;+\; x_N^2
\;-\; 2\sum_{i=1}^{N-1} x_ix_{i+1}\\
&\;=\; x_1^2 \;+\; \sum_{i=1}^{N-1} \left(x_i^2 - 2x_ix_{i+1} + x_{i+1}^2\right) \;+\;
x_N^2\\
&\;=\; x_1^2 \;+\; \sum_{i=1}^{N-1} \left(x_i - x_{i+1}\right)^2 \;+\; x_N^2\\
&\;\geq\; 0, \quad \forall\, x \in \R^N.
\end{align*}

Moreover,
\begin{align*}
x\tp B_Nx = 0 &\iff x_1^2 \;+\; \sum_{i=1}^{N-1} (x_i - x_{i+1})^2 \;+\; x_N^2 \;=\; 0\\
&\iff x_1 = 0;\ x_N = 0;\ x_i = x_{i+1}, \quad \forall\, i = 1:(N-1)\\
&\iff x_i = 0, \quad \forall\, i = 1:N\\
&\iff x = 0.
\end{align*}

We conclude that $x\tp B_Nx > 0$ for all $x \in \R^N$, $x \neq 0$, and therefore that
$B_N$ is a symmetric positive definite matrix; cf. (5.17).
\hfill$\square$
\par\medskip

% Source wording: "symmetric positive matrices" -- reproduced as printed
% (missing "definite"), likewise "Section 5.2.2" is exactly as printed.
Several equivalent ways to characterize symmetric positive matrices are presented
below. A summary of these conditions and comments on their practical applicability can
be found in Section 5.2.2.

\begin{theorem}
(i) A symmetric matrix is symmetric positive semidefinite if and only if all the
eigenvalues of the matrix are greater than or equal to zero.

(ii) A symmetric matrix is symmetric positive definite if and only if all the
eigenvalues of the matrix are strictly greater than zero.
\end{theorem}

\begin{proof}
Let $A$ be a symmetric matrix of size $n$. Recall from Theorem 5.5 that any symmetric
matrix of size $n$ has a full set of $n$ orthogonal eigenvectors. Let $v_1, v_2,
\dots, v_n$ be $n$ orthogonal eigenvectors of $A$ of norm $1$, and let $\lambda_1,
\lambda_2, \dots, \lambda_n$ be the corresponding eigenvalues. If
$Q = \mathrm{col}\,(v_k)_{k=1:n}$ and $\Lambda = \diag(\lambda_k)_{k=1:n}$, the
diagonal form of $A$ is
\begin{equation}
A \;=\; Q\Lambda Q\tp,
\end{equation}
